\clearpage
\section*{2026-09-16: Ocean Parkway's recorded development boundary}
\textit{Agent-drafted research entry.}

The second documentary pass establishes 20,000 square feet of development land
for the 99-unit Ocean Parkway parent. Production currently carries 3,000 square
feet because the same tax-lot identifier described a different parcel in the
reference map. Three GPT-5.6 Sol agents gathered sources for the ten exact-99
parents still unresolved after the Sullivan review. The supervisor read the
decisive Ocean instruments and accepted this allocation in the audit. Nine
parents in that batch remain open. Across the full 877-parent weighting sample,
158 land cases remain unresolved: 119 historical and 39 post-policy. None of
these flags removes a parent from the production panel.

The January 2024 deed, ACRIS 2024022700757001, page 4, defines old Brooklyn
block 7274 lot 30 as a 100-by-30-foot rectangle and old lot 25 as the surrounding
57,000-square-foot parcel. The August 2026 declaration, 2026080700950005,
page 6, defines new lot 30 as a 200-by-100-foot rectangle. Its street offsets
place that rectangle inside old lot 25, outside the former lot 30. The retained
church parcel is 40,000 square feet. The accompanying development agreement,
2026080700950002, pages 3--6, distinguishes the developer's land from the
church's retained land and conveys 64,000 square feet of residential
development rights. These rights permit building floor space; they add no
ground area. The approximately 20,244-square-foot map polygon is a separate
cartographic measurement. The audit adopts the recorded legal rectangle's
20,000 square feet and retains old lot 25 as its archival source parcel.

The framework holds the observed economic-parent boundary fixed across
regimes and measures historical characteristics within it. Later recorded
plans can establish that boundary. Proof that an identical subdivision or
ownership existed at the reference date is not an additional requirement.
The task is to allocate ground and documented common land among the observed
buildings, retained uses, and other phases. A broader zoning lot can include
land contributing development rights without belonging to the residential
development's ground footprint.

Godwin/Kimberly and Beach 29/30 illustrate the remaining questions. The full
Godwin agreement contains the 99-unit job's drawing and depicts retained
commercial land, but its small labels do not legibly identify the separate
65-unit job. The possible missed parent link remains open. Both Beach City
Planning applications identify the three building parcels, whose recorded
areas total 47,366 square feet. They also describe broader zoning lots with
vacant land and illustrative parking layouts. The fifteen earlier two-family
filings have no first permits; their administrative status does not establish
coexistence with the three permitted large buildings. The common-land allocation
and the earlier proposals' disposition remain unresolved.

\textit{Reproduction.} Working changes based on \texttt{a665b53}; run
\texttt{make dof-site-review}. The two-row
\path{parent_site_document_review.csv} records Sullivan and Ocean. The audit
checks each documented reference parcel and area against the archival data and
retains independent filing conflicts. It now reports 719 passing parents and
237 candidate area corrections. Ocean is the only parent whose existing audit
fields changed during this second pass. Seven saved production-panel and
calibration hashes are unchanged. Original scans, API responses, source dates,
page references, and checksums are preserved in the acquisition task's dated
\texttt{site\_research} folders. The supervisor report records the adopted
findings and the outstanding source for each remaining exact-99 parent.
